\section{Ist-Zustandsanalyse der Anwendung Mailbroadcast}
\label{sec:ist_zustand}
Ziel dieses Kapitels ist es, 

Die Analyse des Weiterentwicklungsbedarfs umfasst primär die tiefgründige Untersuchung des Ist-Zustands der ... Anwendung Mailbroadcast, die ...

Schulung erwähnen

In 2.2 aus der Perspektive der Benutzer kurz vorgestellt hier Ist-Zustand umfassend analysiert


\subsection{Überblick und Einsatzbereiche}
Wie in Abschnitt \ref{subsubsec:mbc_allianz} bereits dargelegt, bieten die in Outlook vorhandenen Verteiler keine ausreichende
Flexibilität in der Gruppenkommunikation, da ein Empfängerkreis ausschließlich über die \gls{oe} der Mitarbeiter und dem
Standort, an dem dieser arbeitet, ermittelt werden kann.

ausschließlich die Ermittlung der Empfänger innerhalb einer
\gls{oe} ermöglichen


\cite{mbcSchulung}

\subsection{Architektur}
\subsection{Bereitgestellte Personaldaten}
\subsection{Benutzeroberfläche in Microsoft Access}
\subsection{Bewertung der vorhandenen Dokumentation}
Weiterentwicklung umfasst auch eine gute Doku
\subsection{Zusammenfassung und identifizierte Schwachstellen}
Nur Daten behalten, die wir brauchen --> Daten von Ausgeschiedenen löschen; Index setzen zur Beschleunigung
Verzögerte Mailzustellung

Auf die Erfüllung der Anforderungen eingehen

bsp mit student in einem Referat aber eigentlich in der Berufsausbildung --> Problem bereits in Strukturverteilern in Outlook vorhanden
aber in MBC nicht explizit adressiert